
%\swapnumbers %such that it is « 1.4 Theorem » instead of « Theorem 1.4 »

% ===========================================
% -------- RACCOURCIS SYMBOLIQUES -----------
% ===========================================

% ssi, ssi de définition et ssi idée
\newcommand{\ssi}{\ensuremath{\ \text{ssi}\ }}
\newcommand{\ssid}{\ensuremath{\overset{\triangle}{\ \text{ssi}\ }}}
\newcommand{\ssiidea}{\ensuremath{\overset{\star}{\ \text{ssi}\ }}}


% i.e. et i.e. de définition
\newcommand{\ied}{\ensuremath{\overset{\triangle}{\ \text{i.e.}\ }}}
\newcommand{\idest}{\ensuremath{\ \text{i.e.}\ }}

% indépendance
\newcommand{\indep}{\mathrel{\perp\!\!\!\perp}}

% insérer facilement une jolie fonction
\newcommand{\nicefunc}[5]{%
  \setlength{\arraycolsep}{3pt}%
  \begin{array}{rcl}
    #1 : #2 & \longrightarrow & #3 \\[3pt]
    #4 & \longmapsto & #5
  \end{array}%
}








% ===========================================
% ----- DÉFINITION DES THÉORÈMES ETC. -------
% ===========================================
\usepackage{amsthm}

\theoremstyle{definition}% use "plain" for default
\newtheorem*{thm}{Théorème}
\newtheorem*{lem}{Lemme}
\newtheorem*{prop}{Proposition}
\newtheorem*{cor}{Corollaire}
\newtheorem*{KL}{Klein’s Lemma}
\newtheorem*{conjecture}{Conjecture}

\theoremstyle{definition}
\newtheorem*{defn}{Définition}
\newtheorem*{exmp}{Exemple}
\newtheorem*{xca}{Exercice}

\theoremstyle{remark}
\newtheorem*{rem}{Remark}
\newtheorem*{note}{Note}
\newtheorem*{case}{Case}

\usepackage[framemethod=TikZ]{mdframed}
\usepackage{xcolor}





\definecolor{lightgray}{gray}{0.95}

\mdfdefinestyle{graybox}{
  backgroundcolor=lightgray,
  linecolor=gray,
  linewidth=0pt, %0.6 for normal line
  roundcorner=0pt,%6 for slightly round
  innertopmargin=0pt,
  innerbottommargin=10pt,
  innerleftmargin=10pt,
  innerrightmargin=10pt
}

\definecolor{lightorange}{RGB}{255, 245, 225}

\mdfdefinestyle{orangebox}{
  backgroundcolor=lightorange,
  linecolor=orange!60!black,
  linewidth=0pt, %0.6 for normal line
  roundcorner=0pt,%6 for slightly round
  innertopmargin=0pt,
  innerbottommargin=10pt,
  innerleftmargin=10pt,
  innerrightmargin=10pt
}


% Définition de la couleur vert menthe clair
\definecolor{lightmint}{RGB}{220, 255, 220}

% Définition du style de boîte vert menthe
\mdfdefinestyle{mintbox}{
  backgroundcolor=lightmint, % Couleur de fond vert menthe clair
  linecolor=green!60!black, % Couleur de la bordure
  linewidth=0pt, % Largeur de la bordure (0pt pour aucune bordure visible)
  roundcorner=0pt, % Coins carrés (0pt) ou arrondis (par exemple 6pt)
  innertopmargin=0pt, % Marge intérieure supérieure
  innerbottommargin=10pt, % Marge intérieure inférieure
  innerleftmargin=10pt, % Marge intérieure gauche
  innerrightmargin=10pt % Marge intérieure droite
}


% Couleur rose très clair (à ajuster si besoin)
\definecolor{lightpink}{RGB}{255, 238, 243}
% Style de boîte rose clair
\mdfdefinestyle{pinkbox}{
  backgroundcolor=lightpink, % Fond rose très clair
  linecolor=pink!50!black, % Bordure rose/gris discrète
  linewidth=0pt, % Pas de bordure visible (mettre 1pt si tu en veux une légère)
  roundcorner=0pt, % Coins carrés (mettre 6pt pour arrondis)
  innertopmargin=0pt,
  innerbottommargin=10pt,
  innerleftmargin=10pt,
  innerrightmargin=10pt
}










\usepackage{etoolbox} % nécessaire pour surroundwithmdframed

\surroundwithmdframed[style=orangebox]{thm}
\surroundwithmdframed[style=graybox]{defn}
\surroundwithmdframed[style=mintbox]{prop}
\surroundwithmdframed[style=mintbox]{cor}
\surroundwithmdframed[style=pinkbox]{conjecture}




% ===========================================
% ---------- BOÎTE DE DÉMONSTRATION ----------
% ===========================================

% Définition du style de boîte pour les démonstrations
\mdfdefinestyle{dembox}{
  backgroundcolor=none,      % fond transparent
  linecolor=black,           % couleur de la ligne
  linewidth=3pt,             % épaisseur de la ligne
  leftline=true,             % ligne uniquement à gauche
  topline=false,             % pas de ligne en haut
  bottomline=false,          % pas de ligne en bas
  rightline=false,           % pas de ligne à droite
  innertopmargin=0pt,        % marges internes
  innerbottommargin=10pt,
  innerleftmargin=10pt,
  innerrightmargin=10pt
}

% Définition de l'environnement "dem" (pour démonstration)
\newenvironment{dem}[1][Démonstration]{%
  \begin{mdframed}[style=dembox]
  \noindent\textbf{#1.}\quad
}{%
  \hfill $\square$
  \end{mdframed}
}


